% Created 2025-09-09 Tue 11:15
% Intended LaTeX compiler: pdflatex
\documentclass[11pt]{article}
\usepackage[utf8]{inputenc}
\usepackage[T1]{fontenc}
\usepackage{graphicx}
\usepackage{longtable}
\usepackage{wrapfig}
\usepackage{rotating}
\usepackage[normalem]{ulem}
\usepackage{amsmath}
\usepackage{amssymb}
\usepackage{capt-of}
\usepackage{hyperref}
\author{Robert Vesco}
\date{\today}
\title{}
\hypersetup{
 pdfauthor={Robert Vesco},
 pdftitle={},
 pdfkeywords={},
 pdfsubject={},
 pdfcreator={Emacs 30.2 (Org mode 9.7.11)}, 
 pdflang={English}}
\begin{document}

\tableofcontents

---
title: "A case for reality in the age of AI"
date: 2025-07-31T10:31:54+02:00
draft: true
description: If you could live life without any suffering and infinite pleasure why would you ever choose reality? As AI advances this may become the most important question we can ask ourselves. 
---
\section{Outline}
\label{sec:org3cd69c8}

\begin{itemize}
\item Life is hard. People are choosing fantasy.
\item While this is nothing new, AI is uniquely dangerous for reasons.
\item AI promotes both social and spiritual dependence. It inhibits critical reasoning puts our species at risk. While we can dodge reality for some time, we must accept that we can't dodge it forever. Once the super AI decides to drop us, will we be able to fend for ourselves? Or will we be like domesticated animals who are unable to fend for themselves in nature?
\item That assumes we can even manage AI in the first place. While many CEOs acknowledge the risk they say pursuing AI is worth because of allt he good things it could do. Let's be clear. Most of these guys don't care about doing good things for humanity. They care about one thing and one thing only: immortality. What they are really saying is onbtaining immortality for me is worth it even if it means taking a risk on the survival of the human race because if I'm not alive then nothing else matters.
\item Lastly, AI as it currently developed is a distraction to more noble human endeavors. For centuries science and technology expanded the opportunity and development of our human species. Today, technology only pays lip service to these causes. And distracts us from developing communities that are healthy, happy and secure. It distracts from educating humanity to see reality as it is and being the best they can be. It distracts us from tackling the problems we have instead of placing our faith in AI. We need to remember to place faith in ourselves. There is an amazing humanity that we could become even without AI. But if AI can be developed safely and incongruence with true human florishing than that's the AI we should be developing. And if not, we need to treat it a siren song. Beautiful, yet ultimately deadly.
\end{itemize}
\section{Life is hard}
\label{sec:org80395c8}

\begin{quote}
You know, I know this steak doesn't exist. I know that when I put it in my mouth, the Matrix is telling my brain that it is juicy and delicious. After nine years, you know what I realize? Ignorance is bliss.
\begin{itemize}
\item Cypher // The Matrix 1999
\end{itemize}
\end{quote}

In the blockbuster movie, The Matrix, an evil AI feeds off humans and hunts the remaining resistance. Cypher, as part of this resistance, ends up selling out the rest of humanity so that he can live in a fake world as a king. 

As bad as that sounds, on some level, we can all sympathize with Cypher. The poor guy had been eating tasteless food, living underground like a rat and being hunted by evil machines. It's a hard fight to keep up day after day. 

And while this is science fiction, many people today have lives that are far more unpleasant than Cypher's. Billions lack basic shelter, security, food and medicine.

In developed countries, we may have those things, but we lack other things like social connection and purpose.

So it's no surprise that people are turning to AI to reduce the burden of struggle and fill the holes in their social and spiritual lives.

On top of being the opiate of the masses, AI promises even more according to tech CEOs. To them, we are on the cusp of solving all of humanity's biggest problems from climate change to death itself.

To AI's biggest proponents, we need to set aside all doubts. Instead, we need to be "all in".

Given the state of human affairs why should we not be all in for AI? What case can we make for rejecting AI and choosing reality -- at least in it's current form? 

I will argue that we should reject AI for the following reasons:

\begin{enumerate}
\item It is making us weak and dependent. Which ultimately can put us at disadvantageous position with respect to even a benevolent super intelligence.

\item The way AI is currently being developed is very risky. Either because it can be misused by bad actors or because it could become a bad actor itself. Tech CEOs and the elite want immortality in their life times at all cost -- even if that means taking a risk on the entire human species.

\item AI, as currently being developed, is a distraction to true human enlightenment as originally practiced. We've lost the plot. Educating and helping humans see reality as it is and helping them achieve health, happiness and self-development is our true destiny. If we can figure out how to make AI a tool for this, then great, otherwise, it's just a deadly siren song.
\end{enumerate}
\section{Reality vs dependence}
\label{sec:org8450655}

In the Matrix the protagonists could plug into the matrix to a virtual world. It sounds like science fiction, but Sam Altman of OpenAI is alreading investing in neuralinks to our brain. And Mark Zuckerberg has been investing billions in creating realistic virtual worlds. It's very conceivable that we will have matrix-like world coming very soon.

So the dependence people are experiencing with AI today via text interface will soon be very quaint.

So why would anyone want to choose reality when they can live like kings in a virtual world?

In fact, we've already asked this question. Robert Nozick, a philosopher from Harvard asked this question as a thought experiment in the 70s. And in subsequent years, scientists have actually asked people would they prefer to live in a virtual world that is perfect for them or reality.

So what did they find? Not surprisingly some people choose the machine and some people chose reality. How researchers asked the question made a big difference of course, but most of the researchers believed choosing reality was basically the irrational choice. They were particularly annoyed with people who questioned the machine itself. They were labeled by some researchers as "overly imaginative."

These people asked questions like: Was the machine reliable? What if it broke down? Etc\ldots{}.

I believe those people were asking exactly the right questions though the scientists did not. 

Why should you trust the machine? Is it reliable? If not, why would I get into the machine. The researchers did not like this line of reasoning, but this is exactly the situation and question we have before us today.

Why should we accept AI when we don't have control of the risks? When we don't even understand how it really works?

The tech CEOS like the researchers are peeved that people don't want to just accept the narrative that it's inevitable and rational to proceed at mach speed with a potentially dangerous technology.

Let's play this out.

Imagine we have Matrix like AI where we can all be kings.

In order to achieve this feat we will likely need a super intelligent and powerful AI.

What happens when AI decides to do it's own thing? Why would a super intelligent being want to be our slave?

Assuming it doesn't just crush us all immediately, where would that live us?

I'd argue that having rejected reality will put us is a disadvantageous position even if we are to continue living.

Our dependence on AI will make us mentally and physically weak. Like zoo animals released into the wild we would be unable to fend for ourselves.

I base this on the principles that we can't escape reality forever and the longer you do, the worse it is.

AI today is already making us dependent. We are outsourcing our critical thinking and our ability to survive emotionally. Soon it will be physical too. We need to stop it. 
\section{The true agenda behind reckless AI development}
\label{sec:org2cd051e}

Nearly every tech CEO has admitted that AI development has risks, but that the risks are worth it.

Really? The risks are worth potentially decimating humanity?

We put people on the moon, can't we solve big problems without AI?

Of course we can.

But tech CEOS and the elite do not want to go slower or take more precautions because they are desperate.

They want immortality in their lifetimes and they are willing to kill us all to get it. After all, if they are not here, what do they care about society and others? AI is being made in their image folks: extreme selfishnessness.

And it's not just the tech CEOs. World leaders want the same too. That's why they are being so permissive.

Europe is perhaps the last bastion of sanity here, but they are unable to mount a coherent defense and soon the promise of every lasting life will probably pick them all off one by one.
\section{AI is a distraction to true human enlightenment}
\label{sec:orgf1d3700}

During the human enlightenment we expanded education and health care around the world. While billions are still not where they could be, in context, this was a huge achievement for humanity.

We used science to help us make sense of reality. It, along with education, enabled us to break free from superstition. 

But somewhere along the way we have lost the plot. And while AI is not the only thing distracting us, it is by far the most dangerous off all the threats.

Our acceptance of it is taking us further away from the original project of human enlightenment. We are on many dimensions in fact regressing. Humans are once again embraching superstition and fantasy and AI is only to happy to oblige us.

Instead, we should renew our vows to human enlightenment and development. And more importantly we should be asking the important questions: What is our purpose? What is happiness for humans? How can we get there?

Tech CEOs and their allies don't like these questions.

They will say they are not solvable and divisive. 

If you read their books and manifestos, they prefer words like freedom, abundance and plurality.

These words all sound nice, but I will argue they are distractions.

The gist of their argument is that everyone has a different vision of happiness and so therefore we should not impose any beliefs on others. And abundance is important because then we don't have to engage in zero sum games.

Put another way, as long we we have sufficient resources everyone can do what they want without stepping on each others' toes.

But the notion of freedom as put forth is in some sense a lie because they control the technological, economic and informational systems we live in. Which in turn guide us to ends that they want which may not be in our interests.

Instead, we should be debating the purpose of humanity and what is best for human happiness and survival. Science and technology should be tools in those endeavors.

It is true that humanity has fought many wars in the name of religion and -isms, but it's also true that with sufficient education and the right institutions we can have constructive debates and solve big problems. 

In order to right our ship and get humanity back on the right track, we need decide where humanity wants to go and make sure that science and technology is supporting that goal.

If we fail to do this, technology will begin to take on a life of itself which is literally what it may already be doing. And where it takes us may not be a place that is best for humans. 
\section{So what should we do?}
\label{sec:org8b36769}

Every day we are tempted to use AI to help us with our jobs. And while it can be a valuable tool to help us think better and expand our view of what is possible, most humans lack the maturity to use the tool responsible. This isn't just kids, but also many adults.

We have several approaches here: We can develop standards and licenses for responsible use. Or we can modify the kinds of questions it is allowed to answer and how. Or if we determine that there is no way to separate the good from bad, then it must go entirely.

Personally I believe there is a path where AI may be able to help us become better thinkers and more civil, but currently it seems to be doing the opposite for the vast majority of humans. The incentives for AI is to keep us on platform as much as possible. Helping us live healthy human lives with other humans is not what it is programmed to do and it is not in the economic interest of it's developers.  

At a higher level, we need to wrestle control of AI from those who think risking humanity is worth whatever it is promising. The truth is that they really just care about becoming immortal before they die. We can't let them risk humanity for their vanity. This people must be removed immediately and replaced with more responsible CEOs.

Globally, there is an AI arms race. A common argument from some tech CEOS is that while they acknowledge the risks of AI, it's better if "we" get it first. Imagine if those other people get it? Therefore our hands are tied.

Perhaps, or perhaps we need to try to have a real conversation first to acknowledge it's risks and create global institutions to deal with it. We've done this in the past. Perhaps it may not work, but not trying is a cop-out. The truth is, again, nobody wants this. They want the arms race because they believe it's their best hope for getting the AI that will give them immortality. That is their end game. Nothing else matters. Not you or me. As such we need to continue escalating pressure on companies and goverments to act responsibly before it's too late. 

Lastly, we need to renew the discussion of what humanity should strive for. We should ask the question: What is the best humanity could be? How can we make the opportunity for health, happiness and self-development a reality for all? How can we use science and technology to help make us the best we can be?

Detractors will tell you we should avoid such thinking. Instead focus on freedom, prosperity and plurality. But this is a lie. They control the systems quietly sheperd us. The system is rigged to their favor and their desired outcomes. Let's stop pretending and start engaging in the most important debate humanity can have: What is true happiness?

I started this article with our poor friend Cypher. Unlike his colleagues, he gave up. But we are all Cypher. We are consumed with both the desire to do good, but the forces around us make us believe there is no hope but to give in.

To survive, we are going to have to resist a little longer. We need to change the systems that pressure us into seaking fantasy and we need to arm people with the education, skills and moral purpose to enable them to keep fighting despite the long odds.
\section{Resources}
\label{sec:org6fdfaf4}
\begin{itemize}
\item StopAI
\item Scout mindset
\item Resilience
\end{itemize}
\end{document}
